%% kapa: exam=1 round=34 year=2023 period=1 form=A subject=civil_law subject_ko=민법 seq=28 number=28
%% source: Q-Net 감정평가사 기출문제 (cst00309) artlSeq=5211789 file=raw/1008/5211789/2023년도 제34회 감정평가사 1차시험 1교시.pdf
%% conversion: kordoc 4.24.51 -> tools/gen1.py (auto); review=proofread-sonnet; flags=none
\begin{question}{28}
상린관계에 관한 설명으로 옳지 않은 것은? (다툼이 있으면 판례에 따름)
\begin{choices}
  \item 인접하는 토지를 소유한 자들이 공동비용으로 통상의 경계표를 설치하는 경우, 다른 관습이 없으면 측량비용은 토지의 면적에 비례하여 부담한다.
  \item 지상권자 상호간에도 상린관계에 관한 규정이 준용된다.
  \item 주위토지통행권은 장래의 이용을 위하여 인정될 수 있으므로, 그 범위와 관련하여 장래의 이용상황까지 미리 대비하여 통행로를 정할 수 있다.
  \item 건물을 축조함에는 특별한 관습이 없으면 경계로부터 반미터 이상의 거리를 두어야 한다.
  \item 경계에 설치된 경계표나 담은 특별한 사정이 없는 한, 상린자의 공유로 추정한다.
\end{choices}
\end{question}
\answer{3}
